\section{讨论}
\label{sec:discussion}

我们的发现对未来鲁棒KD方法的设计具有重要启示，同时也突出重要局限性和未来工作方向。

\subsection{方法设计启示}

\textbf{优先转移：从对齐到遮挡处理。}跨域KD研究的主流叙事强调分布对齐——在域偏移下匹配教师和学生输出。我们的结果提示，对于可见度退化，这种关注可能不足。当遮挡移除视觉信息时，对齐方法试图匹配越来越不可靠目标的分布。

\textbf{Phase 1的补充意义：}控制实验通过直接操控遮挡率，为M2假设提供了额外证据。其主要价值不在于重写全文结论，而在于说明：当空间可见信息被进一步遮蔽时，定位蒸馏的相对优势会减弱，这与主实验中的机制判断保持一致。

相反，未来方法应优先\textbf{遮挡感知设计}或\textbf{对空间信息损失的鲁棒性}：
\begin{itemize}
    \item \textbf{遮挡感知蒸馏：}显式建模遮挡概率并相应调整蒸馏损失，对重度遮挡区域降低蒸馏权重。
    \item \textbf{多尺度空间蒸馏：}在多个尺度上分布蒸馏目标，使定位能够利用任何剩余可见信息。
    \item \textbf{部分监督：}而非强制完整特征或输出匹配，允许在可靠教师信号存在的地方进行部分转移。
    \item \textbf{自适应KD权重：}基于遮挡率动态调节$\beta$（KD权重），在低遮挡时充分利用教师知识，高遮挡时更多依赖学生自主学习。
\end{itemize}

\textbf{定位值得更多关注。}我们的结果显示定位蒸馏在轻-中度退化下最鲁棒。这与文献对logit和特征蒸馏的强调形成对比。对于恶劣条件，定位聚焦的KD可能提供更好的基础。

\textbf{退化下的教师质量。}我们的分析揭示一个根本约束：在重度可见度损失下，即使是表现最好的KD分支也无法维持对仅学生基线的增益。这提示教师质量本身显著退化。未来工作可能探索：
\begin{itemize}
    \item \textbf{多教师蒸馏：}结合清晰天气和雾天训练的教师，提供互补信号。
    \item \textbf{教师适应：}在蒸馏前在退化数据上微调或适应教师。
\end{itemize}

\subsection{局限性和说明}

\textbf{机制识别范围。}本文在所测试机制中识别遮挡为最直接得到支持的机制，但未建立完整理论框架。遮挡-定位相关虽强，但相关不蕴含因果，其他混杂因素可能贡献。此外，我们具体测试遮挡作为可见度退化的成分；对比度下降和颜色偏移等其他方面未直接评估，可能贡献于KD失效。

\textbf{Phase 1的作用边界：}通过直接操控遮挡率（而非依赖大气散射模型），Phase 1控制实验部分缓解了“雾天退化来源混杂”的问题，使我们能够更直接地考察遮挡因素本身。但这一补充实验仍应被理解为对M2的做实，而非对主实验主线的替代；真实雾天在密度、颜色和空间分布上仍更复杂，我们识别的机制仍需在真实退化数据上进一步验证。

\textbf{单一数据集和任务。}结果从Foggy Cityscapes（街景数据集）获得。其他领域（如海上、航空）可能表现不同退化模式。我们的发现向其他可见度退化条件（雨、雪、霾）的泛化性需要进一步调查。

\textbf{模型架构特异性。}我们评估YOLOv8变体。虽然YOLO广泛使用，架构差异（如基于Transformer的检测器）可能改变退化下的KD行为。

\subsection{未来方向}

\textbf{遮挡感知KD架构。}开发显式估计遮挡概率并相应调节蒸馏强度的KD方法。重度遮挡区域应接受降低的蒸馏权重以防止误差传播。Phase 2实验正在测试不同$\beta$（KD权重）在各遮挡率下的效果，预期将揭示最优自适应策略。

\textbf{用于鲁棒定位的空间注意力。}调查注意力机制是否能帮助定位蒸馏关注可见物体部分，即使在物体被部分遮挡时也能保持空间准确性。

\textbf{退化下的教师-学生共同训练。}而非从固定的清晰天气教师蒸馏，探索共同训练方法，使教师和学生都学习适应可见度变化。

\textbf{真实世界数据上的机制验证。}在真实雾天数据集上复现我们的分支级分析，以确认合成结果泛化到自然条件。

\subsection{更广阔的视角}

本文贡献于日益增长的认知：KD研究必须超越基准性能，理解方法\emph{何时}和\emph{为什么}成功或失效。我们的机制驱动方法为此类调查提供模板：
\begin{enumerate}
    \item 建立观测结构（分支间差异）；
    \item 形成竞争机制假设；
    \item 设计能区分假设的实验；
    \item 接受负面结果——排除不足的解释本身就是进展。
\end{enumerate}

发现分布错配和不确定性放大提供不足解释——虽反直觉，但与鲁棒机器学习的更广泛教训一致：域偏移通常涉及\emph{信息损失}，不能仅通过校准或对齐补偿。认识这一根本限制应指导未来研究走向更现实的目标和更鲁棒的方法设计。
